\section{Introduction}
\label{sec:introduction}

The Information Field Geometry Theory (IFGT) proposed in this paper is a framework for describing information not as a substance but as a geometric structure. IFGT is not an independent foundational theory; it is positioned as an upper-layer effective theory built upon Vortical Enclosure Dynamics (VED), which describes the generation of physical structure from difference, gradient, flow, and closure.

Information has traditionally been understood primarily as a description of states. When modern physics claims that ``matter is also information,'' this information typically functions as a label attached to the state of a system. The information treated in this paper differs from this usage. Here, information refers to a \emph{structural constraint} arising as a trace left behind when difference does not fully dissipate. Such information has no independent substance; it acts by constraining the paths and distributions of difference chains.

Information is not difference itself. Information appears when difference persists in a distinguishable form and becomes referable in relation to other differences. In this sense, information is a trace produced by the persistence of difference and functions as a scaffold for subsequent difference chains. Information is, therefore, the structure that makes the chaining of differences possible.

This structure is not fully closed; rather, it appears as a fluctuation of difference that cannot be eliminated in the process of approaching closure. In this theory, we call such a structure a \emph{quasi-closure}. A quasi-closure retains differences, produces gradients, and induces flows, thereby forming structure. At the same time, this structure is not separated from physical persistence but interacts with it within the same generative process.

This interaction is not unidirectional but circular. The informational structure as a quasi-closure constrains difference chains, biases their flow, and thereby participates in the formation of new structure. The resulting structure manifests as physical persistence and generates further differences. When these differences in turn remain as quasi-closures, information and matter are generated cyclically, transforming into one another.

Information does not exist statically; it becomes manifest at boundaries where structural change occurs. Such a boundary functions as a front of change and affects the direction and intensity of difference chains. This boundary structure is especially prominent in biological and computational systems, where the interference of past traces with present change is observed as fluctuations of behavior or state.

Such information is concretely implemented in the nervous systems of organisms, the memory of computers, and the internal states of transformer models. These are not information itself but instances in which the structure that retains traces of difference and constrains flow is realized within a physical system. This paper does not address these specific implementations but aims to describe the general structure underlying them.

Thus, the information in this theory is understood as a structure of difference that persists in the process of approaching closure. This structure is described as a geometry that produces the flow and constraint of differences. IFGT provides a framework to describe this informational geometric structure with a minimal set of variables and relations.

\subsection{On the notion of information}
\label{subsec:notion_of_information}

The term \emph{information} is used with considerably different meanings across disciplines, and the position of this paper should be located within that landscape. We here organize four principal notions.

\paragraph{Shannon information.}
In information theory, information is defined as a reduction of probabilistic uncertainty, quantified by $H = -\sum_i p_i \log p_i$. Meaning and weight are intentionally excluded from the formalism. This makes the theory mathematically powerful and domain-independent, but it does not address what information \emph{is}, only how much of it is transmitted.

\paragraph{Physical information (``It from Bit'').}
In the tradition of Wheeler and Landauer, information is treated as a label attached to the state of a physical system, with measurable consequences such as thermodynamic cost of erasure. This formulation is rigorous, but its scope is constrained by the requirement of observability: only quantities that can be measured as physical observables enter the theory. As a consequence, aspects of information such as weight (the degree to which a particular informational structure shapes subsequent dynamics) and semantic interactions below a certain threshold tend to be excluded, not because they are absent, but because they are not directly captured by the available measurement formalism.

\paragraph{Affordance (Gibson).}
In ecological psychology, Gibson~\cite{gibson1986} proposed that information resides not in the environment alone nor in the perceiver alone, but in the relational structure between them. An affordance---the possibility of action offered to an organism by its environment---carries a weight that is determined dynamically by context, rather than as a fixed physical magnitude. This captures an aspect of information that the physical formulation cannot: relationality, context-dependence, and dynamic weighting. However, the notion has remained largely descriptive and has not been given a dynamical formulation.

\paragraph{The informal notion.}
In everyday usage, information is something that persists as meaning and continues to influence subsequent states. This notion is closer to the structural and dynamic aspects that the strict physical formulation leaves out, and it retains the relational character emphasized by affordance.

\paragraph{Position of this paper.}
The information treated in IFGT is not Shannon information, nor a physical state label, but can be understood as a dynamical formulation of the relational structure suggested by affordance and by the informal notion. Concretely, it refers to the \emph{structural constraint that persists as a quasi-closure}---a residue of difference that has not fully closed and that biases subsequent difference chains. The weight of such information is expressed geometrically, as the local depth of the closure-degree field $C$ (see~§\ref{subsec:asymmetry} and~§\ref{subsec:circulatory_structure}). In this sense, IFGT aims to provide a dynamical description of the informational structure whose relational and weighted character has been recognized, but not formalized, in the preceding traditions.
